\documentclass[a4paper]{article}

\newcommand{\docTitle}{Autonomes Fahren}

\usepackage{datetime2}  % for dates
\usepackage[utf8]{inputenc} % UTF-8 input encoding
\usepackage[T1]{fontenc} % fixes \textquotedbl with pdflatex/OT1
\usepackage{textcomp}   % additional text symbols for typewriter/upquote
\usepackage{xcolor}     % for gray color
\usepackage{listings}
\usepackage{marginnote} % for horizontal margin notes
\usepackage{parskip}
\usepackage{amsmath,amssymb,mathtools}
\usepackage{everypage}
\usepackage{fancyhdr}
\usepackage{amsfonts}   % Mengen and natural numbers
\usepackage{enumerate}  % Custom Enumerations
\usepackage[inline]{enumitem}   % List spacing and labels
\usepackage{multicol}   % Multiple Columns
\usepackage{tabularx}

\usepackage{hyperref}   % Links
\hypersetup{
    pdftitle={\docTitle},
    pdfauthor={Moritz Köhler},
    pdfkeywords={DHBW, Hochschule},
    colorlinks=true,
    linkcolor=black,
    filecolor=magenta,
    urlcolor=cyan
}

% -------
% Images
% -------

\usepackage{import}
\usepackage{xifthen}
\usepackage{pdfpages}
\usepackage{tikz}
\usetikzlibrary{arrows.meta, positioning}

\newcommand{\incfig}[1]{%
    \def\svgwidth{\columnwidth}
    \import{./fig/}{#1.pdf_tex}
}

% ---------------------------------
% Vertical Side note (Bottom Left)
% ---------------------------------

\newcommand{\verticaldocinfo}{%
  \rotatebox{90}{\tiny\color{gray}\texttt{\DTMtoday\_\docTitle\_\jobname}}%
}

\AddEverypageHook{
  \put(-40pt,-650pt){\verticaldocinfo}%
}

% ---------------
% Math Shortcuts
% ---------------

% Number sets
\newcommand{\R}{\mathbb{R}}
\newcommand{\N}{\mathbb{N}}
\newcommand{\Z}{\mathbb{Z}}
\newcommand{\Q}{\mathbb{Q}}
\newcommand{\C}{\mathbb{C}}

% Vectors and matrices
\newcommand{\vect}[1]{\mathbf{#1}}
\newcommand{\mat}[1]{\mathbf{#1}}

% Derivatives
\newcommand{\dd}{\mathrm{d}}
\newcommand{\dpartial}[2]{\frac{\partial #1}{\partial #2}}
\newcommand{\ddt}{\frac{\dd}{\dd t}}

% Norms and absolute values
\newcommand{\norm}[1]{\left\lVert #1 \right\rVert}
\newcommand{\abs}[1]{\left\lvert #1 \right\rvert}

% Probability & Statistics
\newcommand{\E}{\mathbb{E}}
\newcommand{\Var}{\mathrm{Var}}
\newcommand{\Prob}{\mathbb{P}}

% Fields and Algebras
\newcommand{\K}{\mathbb{K}}
\newcommand{\F}{\mathbb{F}}

% Set Operations
\newcommand{\compl}[1]{#1^{\mathrm{c}}}

% Matrix and Vector Operations
\newcommand{\T}{^{\top}}
\newcommand{\inv}{^{-1}}
\newcommand{\conj}[1]{\overline{#1}}

% -------------------
% Main lecture macro 
% -------------------

\newcommand{\lecture}[1]{%
  \protect\marginnote{\protect\tiny\protect\color{gray}#1}
}

% -----------------------------
% Command for box with heading
% -----------------------------

\fboxsep2mm
\newcommand{\know}[2]{
    \noindent\fbox{
        \parbox{\textwidth-21pt}{

            \textbf{#1:}
            \vspace{0.125cm}

            #2
        }
    }
}

% -------------------
% Listings Languages
% -------------------

% Shared defaults for code listings
\lstset{
    basicstyle=\ttfamily\small,
    keywordstyle=\color{blue}\bfseries,
    commentstyle=\color{gray},
    stringstyle=\color{teal},
    numberstyle=\tiny\color{gray},
    numbers=left,
    stepnumber=1,
    numbersep=8pt,
    showstringspaces=false,
    frame=single,
    breaklines=true,
    tabsize=4,
    keepspaces=false,
    columns=flexible,
    upquote=true,
    extendedchars=false % disable extended char handling to avoid UTF-8 conflicts
}

% SQL syntax highlighting
\lstdefinelanguage{SQL}{
    morekeywords={SELECT,FROM,WHERE,INSERT,INTO,VALUES,UPDATE,SET,DELETE,CREATE,TABLE,PRIMARY,KEY,NOT,NULL,AND,OR,JOIN,LEFT,RIGHT,INNER,ON,AS,ORDER,BY,GROUP,HAVING,LIMIT},
    sensitive=false,
    morecomment=[l]{--},
    morestring=[b]'
}

% JSON syntax highlighting
\lstdefinelanguage{json}{
    morestring=[b]",
    stringstyle=\color{teal},
    comment=[l]{//},
    morecomment=[s]{/*}{*/},
        morekeywords={true,false,null},
    literate=
     *{0}{{{\color{blue}0}}}{1}
      {1}{{{\color{blue}1}}}{1}
      {2}{{{\color{blue}2}}}{1}
      {3}{{{\color{blue}3}}}{1}
      {4}{{{\color{blue}4}}}{1}
      {5}{{{\color{blue}5}}}{1}
      {6}{{{\color{blue}6}}}{1}
      {7}{{{\color{blue}7}}}{1}
      {8}{{{\color{blue}8}}}{1}
      {9}{{{\color{blue}9}}}{1}
            {:}{{{\color{black}:}}}{1}
            {,}{{{\color{black},}}}{1},
    showstringspaces=false
}

% Language specific styles
\lstdefinestyle{javastyle}{
    language=Java,
    basicstyle=\ttfamily\small,
    keywordstyle=\color{blue}\bfseries,
    commentstyle=\color{gray},
    stringstyle=\color{teal},
    numberstyle=\tiny\color{gray},
    numbers=left,
    stepnumber=1,
    numbersep=8pt,
    showstringspaces=false,
    frame=single,
    breaklines=true,
    tabsize=4,
    keepspaces=true,
    columns=fullflexible,
    upquote=true
}

\lstdefinestyle{sqlstyle}{
    language=SQL,
    basicstyle=\ttfamily\small,
    keywordstyle=\color{blue}\bfseries,
    commentstyle=\color{gray},
    stringstyle=\color{teal},
    numberstyle=\tiny\color{gray},
    numbers=left,
    stepnumber=1,
    numbersep=8pt,
    showstringspaces=false,
    frame=single,
    breaklines=true,
    tabsize=4,
    keepspaces=true,
    columns=fullflexible,
    upquote=true
}

\lstdefinestyle{pythonstyle}{
    language=Python,
    basicstyle=\ttfamily\small,
    keywordstyle=\color{blue}\bfseries,
    commentstyle=\color{gray},
    stringstyle=\color{teal},
    numberstyle=\tiny\color{gray},
    numbers=left,
    stepnumber=1,
    numbersep=8pt,
    showstringspaces=false,
    frame=single,
    breaklines=true,
    tabsize=4,
    keepspaces=true,
    columns=fullflexible,
    upquote=true
}

\lstdefinestyle{cstyle}{
    language=C,
    basicstyle=\ttfamily\small,
    keywordstyle=\color{blue}\bfseries,
    commentstyle=\color{gray},
    stringstyle=\color{teal},
    numberstyle=\tiny\color{gray},
    numbers=left,
    stepnumber=1,
    numbersep=8pt,
    showstringspaces=false,
    frame=single,
    breaklines=true,
    tabsize=4,
    keepspaces=true,
    columns=fullflexible,
    upquote=true
}

\lstdefinestyle{bashstyle}{
    language=bash,
    basicstyle=\ttfamily\small,
    keywordstyle=\color{blue}\bfseries,
    commentstyle=\color{gray},
    stringstyle=\color{teal},
    numberstyle=\tiny\color{gray},
    numbers=left,
    stepnumber=1,
    numbersep=8pt,
    showstringspaces=false,
    frame=single,
    breaklines=true,
    tabsize=4,
    keepspaces=true,
    columns=fullflexible,
    upquote=true
}

\lstdefinestyle{jsonstyle}{
    language=json,
    basicstyle=\ttfamily\small,
    keywordstyle=\color{blue}\bfseries,
    commentstyle=\color{gray},
    stringstyle=\color{teal},
    numberstyle=\tiny\color{gray},
    numbers=left,
    stepnumber=1,
    numbersep=8pt,
    showstringspaces=false,
    frame=single,
    breaklines=true,
    tabsize=4,
    keepspaces=true,
    columns=fullflexible,
    upquote=true
}

% Default listing style
\lstset{language={}}

\title{\docTitle}
\author{Moritz}
\date{\today}

\begin{document}
\maketitle
\tableofcontents

\lecture{2026-10-05}

\section{Bildverarbeitung}

Typische Pipeline:

\begin{enumerate}
    \item Bildaufnahme
    \item Vorverarbeitung
    \item Merkmalsextraktion
    \item (Segmentierung)
    \item (Auswertung)
    \item Klassifikation
\end{enumerate}

\subsection{Bildaufnahme}

Der Mensch hat: Linse, Iris, Stäbchen/ Zapfen

Eine Kamera hat: Objektiv, Blende, Sensor

Die Blende dient dazu um das Licht für den Sensor zu normalisieren (Tag und Nacht soll es ein Objekt im Bild erkennbar machen)

\href{https://www.youtube.com/watch?v=MytCfECfqWc}{The Science of Camera Sensors}

\subsubsection{Lochkameramodell}

\noindent
\begin{minipage}[t]{0.48\textwidth}
    Strahlensatz

    $$\frac{h'}{h} = \frac{f}{d}$$

    Mithilfe von Domain Wissen, wie Schildhöhe bei den Smart Rollers, kann mit dem Strahlensatz auch die Entfernung von Schildern ermittelt werden, obwohl nur eine Kamera verbaut wird.
    
    (x, y) sind Bildkoordinaten

    (u, v) sind Pixelkoordinaten
\end{minipage}
\hfill
\begin{minipage}[t]{0.48\textwidth}
    \vspace{0pt}
    \centering
    \includegraphics[
        page=18,
        width=\textwidth,
        trim=400pt 0pt 0pt 0pt,
        clip
    ]{doc/02_Bildverarbeitung.pdf}
\end{minipage}

\subsubsection{Weitere Formeln/ Begriffe}

Typische Aufgaben/ Fachbegriffe:

\begin{itemize}
    \item Projektion in die Bildebene
    \item Transformation zu Pixelkoordination
    \item Intrinsische Kameramatrix K
    \item Homogene Vektoren
\end{itemize}

\subsubsection{Koordinationsysteme}

\begin{itemize}
    \item Welt
    \item Kamera
    \item Standard Kamera (Achsen sind Horizontal zur Erde)
    \item Straße
    \item Straße (ISO8855) (Zi nach oben, statt nach vorne)
\end{itemize}

Stichworte:

\begin{itemize}
    \item Rotationsmatrix
    \item Transformation
    \item Rotation (roll, pitch, yaw)
\end{itemize}

\subsection{Bildvorverarbeitung}

\subsubsection{Farbräume}

Für die Berechnung von Farbe auf Schwarz weiß werden die RGB Farben nicht jeweils mit 0.33 Multipliziert. Typische Werte sind [0.299, 0.587, 0.144] Dies ist ähnlich zu den Anzahl der Zapfen in unserem Auge.

\subsubsection{Histogramme (Helligkeit)}

Normierung durch die Anzahl der Pixel, damit das Integral 1 ergibt.

% Ich war sehr unaufmerksam. Hier fehlt einiges.

Stichworte:

\begin{itemize}
    \item Thresholding/ Punktoperation
    \item (Lineare) Transformationen
\end{itemize}

\subsection{Merkmalsextraktion}

\subsubsection{Faltung}

Wir können mittels eine Maske ein Bild kleiner machen, indem wir den Mittelwert von der Umgebung in einen neuen Pixel schreiben.

Diese Maske muss nicht zwingend der Mittelwert sein, sondern kann auch explizit gewählt werden um z.B. Kanten zu finden.

\subsubsection{Filter}

Mögliche Filter:

\begin{itemize}
    \item Mittelwert
    \item Gaußscher Glättungsfilter
    \item Schärfungsfilter
    \item Laplacian Kantenfilter
    \item Sobel Kantenfilter
    \item Prewitt Kantenfilter
    \item Rauschunterdrückung
    \item Morphologische Filter (Erosion, Dilatation)
\end{itemize}

Laplacian Kantenfilter (ohne Diagonale):

$$\begin{bmatrix}
    0 & 1 & 0\\
    1 & -4 & 1\\
    0 & 1 & 0\\
\end{bmatrix}$$

Dilation macht Helle Pixel größer/ breiter, damit z.B. ein Kantenfilter besser funktioniert.

Erosion macht Dunkle Pixel größer/ breite.

\end{document}